
\sgaStableTarget{sga6.e10.sec.4.5}\subsection*{4.5}
The equivalence relation defined by the intersection product @@M0@@ is called numerical equivalence:

\sgaStableTarget{sga6.e10.def.4-5-1}\textbf{Definition 4.5.1.} \uline{Two elements} @@M1@@ \uline{and} @@M2@@ \uline{of} @@M3@@ @@M4@@ (resp. \uline{of} @@M5@@\uline{) are said to be numerically equivalent when they define the same homomorphism} @@M6@@ (resp. @@M7@@\uline{). We say that two invertible} @@M8@@\uline{-modules} @@M9@@ and @@M10@@ \uline{are numerically equivalent when their first Chern classes} @@M11@@ \uline{are numerically equivalent.}

\sgaStableTarget{sga6.e10.prop.4-5-2}\textbf{Proposition 4.5.2.} \uline{An invertible} @@M12@@\uline{-module} @@M13@@ \uline{is numerically equivalent to zero if and only if}
\sgaStableTarget{sga6.e10.eq.4-5-3}@@M14@@
\uline{for every} @@M15@@ \uline{and every integral closed curve} @@M16@@.

This follows from 1.1.2 and 4.3.1.

Denote by @@M17@@ the subgroup of @@M18@@ consisting of the classes of invertible @@M19@@-modules numerically equivalent to zero. From 2.4.4 one obtains

\sgaStableTarget{sga6.e10.cor.4-5-4}\textbf{Corollary 4.5.3.} @@M20@@.

In Exp. XIII it will be proved that one even has @@M21@@.

\sgaStableTarget{sga6.e10.sec.5}\section*{5. The isomorphism @@M22@@}

\sgaStableTarget{sga6.e10.sec.5.1}\subsection*{5.1.}
Let @@M23@@ be an arbitrary ringed topos. Then one can define the ring @@M24@@ and the filtration @@M25@@ exactly as for schemes. Likewise, define the group @@M26@@ as the group of isomorphism classes of invertible @@M27@@-modules. Denote by

